\documentclass[
    language=english,
    doctype=short-story,
    institution=none,
]{../../omnilatex}

\RequirePackage{config/document-settings}

\begin{document}
\maketitle

\vspace{1cm}

\begin{center}
\textit{``The truth is rarely pure and never simple.''}\\
--- Oscar Wilde, \textit{The Importance of Being Earnest}
\end{center}

\vspace{1cm}

\section*{THE CARTOGRAPHER'S DAUGHTER}

\noindent By Eleanor Vance

\vspace{0.5cm}

The map arrived on a Tuesday, which seemed wrong. Maps should arrive on mornings---bright, purposeful mornings when the light was clear enough to read by and the day ahead felt like something you could chart. But this map came on a Tuesday afternoon in November, carried by a mail carrier who looked like he'd rather be anywhere else, and it landed on the doormat of Apartment 4B with the quiet finality of a verdict.

Clara found it when she came home from work. She set her bag down, hung her coat on the hook by the door, and noticed the envelope lying at an angle on the mat, as though it had been dropped in haste. It was large, cream-colored, and addressed in handwriting she recognized immediately.

Her father's handwriting.

She hadn't spoken to her father in seven years. Not since the argument at Thanksgiving, the one that had started about something trivial---whether the potatoes were overcooked, or whether the turkey needed another hour, or whether the wine was too warm---and had spiraled into something vast and dark and irreversible. He had said things. She had said things. The silence that followed had calcified into something that felt permanent.

She picked up the envelope. Her name was written on the front in his precise, angular script: \textit{Clara Whitfield}. No ``Dear,'' no ``Ms.,'' just her name, as though it were a coordinate on a grid and he was marking her location.

She carried the envelope to the kitchen, made herself a cup of tea, and sat down at the table. The tea grew cold. She stared at the envelope for twenty minutes before she opened it.

Inside was a map.

Not a printed map, not a road atlas or a subway diagram or a topographical survey. It was a hand-drawn map, rendered in ink on heavy paper, with the meticulous attention to detail that had made her father one of the most respected cartographers in the Pacific Northwest. Every line was precise. Every label was lettered in his careful hand. The scale was exact, the compass rose was ornate, and in the bottom right corner, in small letters, was the date: November 14, 2024.

The map showed a place Clara had never seen. It was labeled \textit{Whitfield} in the center, and it depicted what appeared to be a small town nestled in a valley between two mountain ranges. There were streets, parks, a river, a harbor. There were landmarks---a library, a church, a school, a cemetery. There were paths marked in dashed lines leading into the surrounding hills.

And there was a single red dot, placed at the intersection of two streets in the center of town. The streets were labeled \textit{Elm} and \textit{Fourth}. Beside the red dot, in tiny letters, was written: \textit{Start here}.

Clara turned the map over. On the back was a note, written in the same angular hand:

\begin{quote}
\textit{Clara,}\\[0.3cm]
\textit{I know you have no reason to trust me. I know I have given you no reason to believe that I am anything other than what you think I am.}\\[0.3cm]
\textit{But I need you to come. I need you to follow the map. I need you to see what I've made.}\\[0.3cm]
\textit{There are things I should have told you. Things I couldn't tell you then. Things I can only show you now.}\\[0.3cm]
\textit{Please.}\\[0.3cm]
\textit{--- Dad}
\end{quote}

Clara set the map down and looked out the window. The November light was fading, the sky turning the color of bruises. Somewhere in the distance, a train whistle sounded, long and mournful.

She should throw it away. She should crumple it up, toss it in the trash, and pour the cold tea down the sink. She should go on with her life, her quiet, orderly, predictable life, and forget that her father existed.

Instead, she opened her laptop and searched for \textit{Whitfield, Washington}.

\section*{}

There was no such place. She searched every database she could think of, every mapping service, every geographic registry. There was a Whitfield County in Georgia, a Whitfield Township in Pennsylvania, a Whitfield Street in several cities. But no Whitfield, Washington. No town in a valley between two mountain ranges, with a river and a harbor and a cemetery and a red dot at the intersection of Elm and Fourth.

She looked at the map again. Her father's cartography had always been precise, almost obsessive in its accuracy. He had taught her to read maps when she was seven, had shown her how to use a compass, how to measure distance with a piece of string, how to read contour lines and understand the shape of the land. ``A map is a promise,'' he used to say. ``It tells you where you are and where you can go. It doesn't lie.''

But this map depicted a place that didn't exist. Or---she felt a strange tingling at the back of her neck---a place that didn't exist \textit{yet}.

She printed the map, folded it carefully, and packed a bag.

\section*{}

The drive to Washington took three days. Clara drove north through Oregon, following the highways that wound through mountain passes and along river valleys, past logging towns and reservation lands and the vast, dark stretches of national forest where the trees grew so thick that daylight barely touched the ground.

She followed the map, or tried to. The roads it depicted didn't correspond to any roads on her GPS. She ended up on back roads, dirt tracks, logging roads that the map showed as solid lines but that her car could barely manage. Twice she had to turn back and find alternate routes.

On the third day, she reached a point where the map showed a road leading into a valley. The road wasn't on her GPS. There was no sign, no marker, nothing to indicate that a road existed here at all. Just trees, dense and dark, and the sound of a river somewhere in the distance.

She turned off the highway and drove into the trees.

The road was narrow, barely wide enough for one car, and it wound through the forest in tight switchbacks that made her stomach lurch. The trees pressed in on both sides, their branches interlocking overhead to form a canopy that blocked out the sky. The light was dim, greenish, filtered through layers of leaves and needles.

After twenty minutes, the trees opened up, and Clara found herself looking down into a valley.

It was exactly as the map had depicted. The town of Whitfield lay below her, nestled between two ranges of mountains, with a river running through its center and a harbor at the river's mouth. She could see the streets, the houses, the library with its copper roof, the church with its white steeple. She could see Elm Street, running north-south, and Fourth Street, running east-west, and at their intersection, a small park with a fountain.

There was no red dot, of course. That was just ink on paper. But the place was real. It existed. Her father had drawn a map of a place that existed, and he had sent it to her, and now she was here.

She drove down into the valley.

\section*{}

Whitfield was a town frozen in time. The architecture was early twentieth century---Victorian houses with wraparound porches, a main street with brick storefronts, a movie theater with an art deco marquee. But everything was immaculate, as though the town had been preserved in amber. There were no weeds in the cracks of the sidewalks. No paint peeling from the houses. No trash on the streets.

There were also no people.

Clara drove slowly through the empty streets, her headlights cutting through the early evening gloom. She passed the library, the church, the school. All closed. All dark. All perfect.

She found Elm and Fourth. The intersection was exactly as the map had depicted, down to the park with its fountain. The fountain was running, water arcing gracefully into the air, catching the last light of the day. But the park was empty. No children playing, no couples walking, no old men feeding pigeons.

She parked the car and walked to the fountain. The red dot on the map had marked this spot, this exact spot, as the starting point. She stood where the dot had been and looked around.

In the center of the fountain, beneath the water, something glinted. Clara knelt, reached into the cold water, and pulled out a small metal box, tarnished with age but still sealed. There was a name engraved on the lid: \textit{C. Whitfield}.

She opened the box. Inside was a key, a folded piece of paper, and a small brass compass. The paper was a note, written in her father's handwriting:

\begin{quote}
\textit{The key opens the house at 42 Elm Street. Go there. Read what's inside. Then decide.}\\[0.3cm]
\textit{I'll be waiting.}
\end{quote}

Clara looked at the compass. The needle spun wildly for a moment, then steadied, pointing not north but toward a house on Elm Street---a small Victorian with a blue door and a garden that, despite the November cold, was full of flowers.

She walked to the house, unlocked the door, and stepped inside.

The house was warm. A fire was burning in the fireplace. The furniture was familiar---the same furniture from her childhood home, the same bookshelves, the same photographs on the walls. But the photographs were different. They showed a family---her father, her mother, and Clara---but in settings she'd never seen. They were smiling. They were happy. They were together.

On the mantelpiece was a framed document, an old newspaper clipping with the headline: \textit{Cartographer Completes Life's Work: Entire Town Built from Memory}.

Below the headline was a photograph of her father, younger, standing in front of the fountain at Elm and Fourth, smiling at the camera.

Clara sat down in the chair by the fire and read the article. Her father, it seemed, had spent thirty years building Whitfield---every house, every street, every detail---based on a town he had visited once as a child, a town that had been destroyed by a flood in 1952. He had recreated it from memory, from old photographs, from the stories his own parents had told him. He had built it as a monument to something lost, a place that existed only in memory, made real through cartography and love and an obsessive attention to detail.

And then, the article said, he had invited his daughter to see it.

Clara looked at the photograph on the mantelpiece, at her father's smiling face, and felt something shift inside her, something that had been locked in place for seven years beginning to loosen.

The front door opened.

Her father stood in the doorway, older than she remembered, thinner, his hair completely white now. He looked at her with an expression that was equal parts hope and fear.

``You came,'' he said.

``I came,'' she said.

They looked at each other across the room, across seven years of silence, across the vast and terrible distance that had grown between them.

Then Clara stood up and walked to him, and he opened his arms, and she let herself be held, and outside, the fountain at Elm and Fourth continued to flow, and the town of Whitfield settled into the evening, whole and perfect and waiting.

\end{document}
